\section{TDCA}
\label{sec:method}

\subsection{Dynamic reasoning hypergraph}
We represent state as
\begin{equation}
H_t=(V_t,E_t,\mathcal{B}_t,\mathcal{L}_t,\mathbf{B}_t),
\end{equation}
where $V_t$ contains typed nodes (\textsc{Subgoal}, \textsc{Evidence}, \textsc{Claim}, and \textsc{Answer}), $E_t$ contains directed hyperedges, $\mathcal{B}_t$ stores active branches, and $\mathcal{L}_t$ is an append-only audit ledger. A hyperedge $e:S_e\rightarrow v_e$ requires a set of premises $S_e\subseteq V_t$, allowing conjunctions that ordinary pairwise graphs obscure. Query decomposition initializes a root objective and two to four explicit subgoals; later operations can add evidence, instantiate claims, merge duplicates, or invalidate dependencies.

Each node $v$ maintains a belief vector rather than one confidence:
\begin{equation}
\mathbf{b}_v=\big[s_v,w_v,h_v,g_v,k_v,d_v,u_v,z_v\big],
\end{equation}
where $s$ is absolute support, $w$ relative weight among alternatives, $h$ entropy, $g$ evidence gap, $k$ contradiction mass, $d$ downstream impact, $u$ dependency unlock, and $z$ propagated heat. Separating these channels prevents, for example, a high relative weight among weak candidates from masquerading as strong absolute support.

\begin{figure}[t]
\centering
\placeholderbox{3.0cm}{\method{} overview: query decomposition $\rightarrow$ dynamic reasoning hypergraph $\rightarrow$ typed diffusion $\rightarrow$ operation--fidelity EVC selection $\rightarrow$ controller-mediated mutation $\rightarrow$ safe terminal decision.}
\caption{Conceptual overview of graph-state-driven computation.}
\label{fig:overview}
\end{figure}

\subsection{Typed directional diffusion}
Local scores alone do not reveal which uncertainty blocks the answer. We propagate selected channels over type-compatible incidence operators. Let $x_t^{(c)}\in\mathbb{R}^{|V_t|}$ be channel $c$ and $P_t^{(c)}$ a row-normalized, directional transition matrix induced by permitted hyperedge roles. Starting from local signal $x_{t,0}^{(c)}$, diffusion uses restart $\rho$ and decay $\gamma$:
\begin{equation}
x_{t,\ell+1}^{(c)}=\rho x_{t,0}^{(c)} +(1-\rho)\gamma P_t^{(c)}x_{t,\ell}^{(c)},\qquad \ell=0,\ldots,L-1.
\label{eq:diffusion}
\end{equation}
The development configuration uses $\rho=0.4$, $\gamma=0.65$, and $L=3$. Direction and type masks prevent indiscriminate smoothing: evidence support may flow to a claim it entails, while answer-level uncertainty flows backward to prerequisite subgoals as computational demand. Diffusion is recomputed after structural events rather than after every scalar update.

\subsection{Explicit joins and event-triggered revision}
A \textsc{Join} operator proposes a new claim only from two to four addressable premises. Join schemas cover path composition, shared-role constraints, set intersection, comparison, and numeric aggregation. A proposal records premise IDs, provenance spans, the transformation schema, and an independently scored raw output. Acceptance requires premise readiness, type compatibility, novelty, support, and absence of unresolved contradiction. Thus an auditable chain has the form
\begin{equation}
\{c_1,c_2,\ldots,c_m\}\xRightarrow[\text{schema},\,\text{provenance}]{\textsc{Join}}c_{m+1},\qquad 2\leq m\leq4.
\end{equation}

The graph editor is event-triggered. Contradiction, belief collapse, changed best candidate, stale dependency, or a failed high-value operation can emit an edit event. Only the controller may commit mutations; worker components return proposals. \textsc{Revise} can retract a derived claim, down-weight a hyperedge, reopen a subgoal, or fork an alternative branch. Every mutation stores its cause and before/after state in $\mathcal{L}_t$.

\subsection{Multi-resource expected value of computation}
At each step, the controller constructs feasible packets $p=(o,f,r)$ consisting of an operation $o$ (retrieve, verify, join, revise, or answer-related checks), a fidelity $f$, and a target graph region $r$. Each packet has normalized benefit components and separately normalized costs. We use a deterministic additive score:
\begin{align}
\widehat{\mathrm{EVC}}_t(p)={}&\alpha_1\widetilde{\Delta\mathrm{support}}+\alpha_2\widetilde{\Delta\mathrm{gap}}+\alpha_3\widetilde{\Delta\mathrm{entropy}}+\alpha_4\widetilde{\mathrm{unlock}}+\alpha_5\widetilde{\mathrm{downstream}} \nonumber\\
&-\lambda_{\mathrm{call}}\widetilde c_{\mathrm{call}}-\lambda_{\mathrm{tok}}\widetilde c_{\mathrm{tok}}-\lambda_{\mathrm{ret}}\widetilde c_{\mathrm{ret}}-\lambda_{\mathrm{risk}}\widetilde c_{\mathrm{risk}}.
\label{eq:evc}
\end{align}
All components are normalized before addition; the distinct shadow prices $\lambda$ preserve the resource accounting. Priors are operation-, fidelity-, and context-specific. After execution, actual cost and realized graph utility are logged, and a bounded within-question update adjusts later estimates. No model parameters are trained, and feedback does not cross evaluation examples.

\begin{figure}[t]
\centering
\fbox{\begin{minipage}{0.94\linewidth}
\textbf{Algorithm 1: Graph-state-driven computation allocation}\smallskip

\textbf{Input:} question $q$, corpus $\mathcal{D}$, model $M$, budget $\mathbf{B}$.\\
1. Initialize $H_0$ by decomposing $q$; validate graph invariants.\\
2. For $t=0,1,\ldots$: recompute typed diffusion after structural events.\\
3. If the answer gate passes, return \textsc{Answer}; if no useful packet remains, return \textsc{Abstain}.\\
4. Enumerate feasible operation--fidelity packets under residual budget.\\
5. Predict component-wise benefit, cost, and $\widehat{\mathrm{EVC}}_t(p)$; select the deterministic maximum with stable tie-breaking.\\
6. Execute the packet as a proposal; independently score raw outputs.\\
7. Controller validates and commits accepted mutations, records predicted EVC, actual cost, realized utility, and updates within-question estimates.\\
8. If no packet fits the residual budget, return \textsc{BudgetExhausted}; otherwise continue.
\end{minipage}}
\caption{High-level \method{} inference loop. Detailed schemas, gates, and invariants are in Appendix~\ref{app:details}.}
\label{alg:tdca}
\end{figure}

\subsection{Safe termination}
An answer is emitted only when its candidate exists, required dependencies are complete, absolute support exceeds a threshold, the evidence gap is closed, and no blocking contradiction remains. Formally,
\begin{equation}
\mathrm{AnswerReady}(a)=C_a\land D_a\land S_a\land G_a\land\neg K_a.
\end{equation}
If the graph is coherent but no answer is sufficiently supported and no positive-EVC packet remains, the system abstains. If a potentially useful packet exists but cannot fit the remaining budget, it reports budget exhaustion. These outcomes are evaluated separately rather than collapsing all non-answers into failure.
